% ===========================================================================
%  Chapitre 14 — Statistique à deux variables
%  PE 2026, composante TRAITEMENT DES DONNÉES ET PROBABILITÉS
% ===========================================================================
\bmtchapitre{Statistique à deux variables}
  {Représenter une série statistique double par un nuage de points,
   déterminer le point moyen, ajuster le nuage par une droite — méthode
   graphique ou méthode de Mayer — et estimer une valeur future.}
  {Traitement des données \textbullet\ environ 8 heures}

Deux colonnes de chiffres : l'année et la production, le rang du mois et le
bénéfice, la note de natation et la note de danse. La question est toujours
la même — ces deux grandeurs évoluent-elles ensemble, et si oui, que
prévoit-on pour demain ?

La réponse tient en une droite. On place les points, on les résume par un
point moyen, on trace la droite qui suit le nuage, et l'on prolonge. C'est
le dernier chapitre du programme, et c'est aussi celui dont l'usage dépasse
le plus l'école : un commerçant, un agronome, un directeur d'établissement
font exactement ce calcul.

% ---------------------------------------------------------------------------
\begin{revision}
Le chapitre 2 et quelques acquis de seconde.

\begin{enumerate}[label=\textbf{\arabic*.}, leftmargin=*, itemsep=0.35em]
  \item Calculer la moyenne de $12$, $15$, $18$ et $19$.
  \item Déterminer l'équation de la droite passant par $A(1\,;\,5)$ et
        $B(3\,;\,11)$.
  \item Placer les points $(1\,;\,340)$ et $(2\,;\,310)$ dans un repère où
        $1$ cm représente une unité en abscisse et $10$ unités en ordonnée,
        l'origine des ordonnées étant placée à $300$.
  \item Résoudre $\begin{cases} 2a+b = 325\\ 5a+b = 380\end{cases}$.
  \item Calculer $\dfrac{55}{3}$ à $10^{-2}$ près.
  \item Une production passe de $200$ à $250$ tonnes : de quel pourcentage
        a-t-elle augmenté ?
\end{enumerate}
\end{revision}

\begin{automatismes}
Sans calculatrice.

\begin{enumerate}[label=\textbf{\alph*)}, leftmargin=*, itemsep=0.25em]
  \item Moyenne de $1$, $2$, $3$, $4$, $5$ et $6$
  \item Moyenne de $10$, $20$ et $30$
  \item Coefficient directeur de la droite passant par $(2\,;\,4)$ et
        $(5\,;\,13)$
  \item $\dfrac{380-325}{5-2}$
  \item Augmenter $200$ de $10\,\%$
  \item Diminuer $50$ de $20\,\%$
  \item Coefficient multiplicateur d'une hausse de $25\,\%$
  \item Deux hausses successives de $10\,\%$ : coefficient global
\end{enumerate}
\end{automatismes}

\begin{corrige}[automatismes]
a) $3{,}5$ \quad b) $20$ \quad c) $3$ \quad d) $\frac{55}{3} \approx 18{,}33$
\quad e) $220$ \quad f) $40$ \quad g) $1{,}25$ \quad h) $1{,}21$.
\end{corrige}

% ===========================================================================
\section{Série double et nuage de points}

\begin{definition}[série statistique à deux variables]
On observe, sur une même population, deux caractères numériques $x$ et $y$.
La liste des couples $\left(x_{i}\,;\,y_{i}\right)$ s'appelle une
\textbf{série statistique double}. Représentés dans un repère, ces couples
forment le \textbf{nuage de points} de la série.
\end{definition}

\begin{definition}[point moyen]
Le \textbf{point moyen} du nuage est le point
$\pmoyen\left(\moy{x}\,;\,\moy{y}\right)$, où $\moy{x}$ et $\moy{y}$
désignent les moyennes respectives des $x_{i}$ et des $y_{i}$ :
\[
  \moy{x} = \frac{x_{1}+x_{2}+\cdots+x_{n}}{n}, \qquad
  \moy{y} = \frac{y_{1}+y_{2}+\cdots+y_{n}}{n}.
\]
\end{definition}

\begin{avertissement}
Les sujets imposent l'échelle du tracé et, souvent, une origine décalée :
« sur l'axe des ordonnées, placer $300$ à l'origine puis $1$ cm pour
$10$ élèves ». Cette consigne se respecte à la lettre — elle vaut un point —
et elle se signale sur le dessin par une double barre sur l'axe, juste
au-dessus de l'origine.
\end{avertissement}

\begin{visualisation}[un nuage, son point moyen, sa droite]
\begin{center}
\begin{tikzpicture}
\begin{axis}[
  width = 11cm, height = 6.4cm,
  axis lines = left, xlabel = {rang de l'année $x_{i}$},
  ylabel = {effectif $y_{i}$},
  xmin = 0, xmax = 7.4, ymin = 295, ymax = 410,
  xtick = {1,2,3,4,5,6,7}, ytick = {300,320,340,360,380,400},
  axis line style = {bmtGris},
  tick label style = {font=\scriptsize, color=bmtGris},
  grid = both, grid style = {bmtGrisClair, line width=0.3pt}, clip = true,
]
  \addplot[only marks, mark=*, mark size=1.9pt, bmtIndigo] coordinates {
    (1,340) (2,310) (3,325) (4,380) (5,365) (6,395)};
  \addplot[bmtRaphia, line width=1.2pt, domain=0.6:7.2]
    {18.333*x+288.333};
  \filldraw[bmtLaterite] (axis cs:2,325) circle (2.2pt);
  \filldraw[bmtLaterite] (axis cs:5,380) circle (2.2pt);
  \filldraw[bmtVetiver] (axis cs:3.5,352.5) circle (2.4pt);
  \node[bmtLaterite, font=\scriptsize, anchor=north west]
    at (axis cs:2.1,323) {$G_{1}$};
  \node[bmtLaterite, font=\scriptsize, anchor=north west]
    at (axis cs:5.1,378) {$G_{2}$};
  \node[bmtVetiver, font=\scriptsize, anchor=south east]
    at (axis cs:3.45,354) {$G$};
\end{axis}
\end{tikzpicture}
\end{center}

Les six points bleus forment le nuage. Le point vert $G$ en est le point
moyen : il se trouve toujours « au centre » du nuage. Les deux points rouges
$G_{1}$ et $G_{2}$ sont les points moyens des deux moitiés ; la droite qui
les joint est la droite d'ajustement de Mayer, et elle passe elle-même par
$G$.
\end{visualisation}

% ===========================================================================
\section{Ajuster le nuage par une droite}

\begin{methode}{la méthode de Mayer, en cinq gestes}
\begin{enumerate}[label=\textbf{\arabic*.}, leftmargin=*, itemsep=0.15em]
  \item \emph{Partager} la série en deux sous-groupes de même effectif : la
        première moitié des valeurs, puis la seconde.
  \item \emph{Calculer} les points moyens $G_{1}$ et $G_{2}$ de chaque
        sous-groupe.
  \item \emph{Placer} $G_{1}$ et $G_{2}$ sur le graphique et tracer la
        droite $\droiteMayer$.
  \item \emph{Déterminer son équation} $y = ax+b$ en résolvant le système
        formé par les deux points — c'est exactement un système
        $2\times2$ du chapitre 2.
  \item \emph{Estimer} : pour une valeur future de $x$, calculer $ax+b$ ;
        pour une valeur de $y$, résoudre l'équation en $x$.
\end{enumerate}
\end{methode}

\begin{propriete}[la droite de Mayer passe par le point moyen]
Lorsque les deux sous-groupes ont le même effectif, le point moyen
$\pmoyen$ de la série entière appartient à la droite $\droiteMayer$.
\end{propriete}

\begin{demonstration}
Les deux sous-groupes ayant $\frac n2$ éléments chacun, la moyenne générale
est la moyenne des deux moyennes :
$\moy{x} = \frac{x_{G_{1}}+x_{G_{2}}}{2}$ et de même pour $y$. Le point
$\pmoyen$ est donc le milieu du segment $\left[G_{1}G_{2}\right]$, et un
milieu appartient toujours au segment.
\end{demonstration}

\begin{remarque}
Cette propriété fournit une vérification gratuite : après avoir trouvé
l'équation, on remplace $x$ par $\moy{x}$ et l'on doit retrouver $\moy{y}$.
Dix secondes, et l'on sait si l'on s'est trompé.
\end{remarque}

\begin{exemple}[une série complète, du tableau à la prévision]
On reprend les effectifs du nuage précédent : rangs $1$ à $6$, effectifs
$340$, $310$, $325$, $380$, $365$, $395$.

\emph{Point moyen.} $\moy{x} = \frac{1+2+3+4+5+6}{6} = 3{,}5$ et
$\moy{y} = \frac{2\,115}{6} = 352{,}5$, donc
$\pmoyen(3{,}5\,;\,352{,}5)$.

\emph{Sous-groupes.} $G_{1}$ résume les trois premières années :
$\left(2\,;\,\frac{975}{3}\right) = (2\,;\,325)$. $G_{2}$ résume les trois
dernières : $\left(5\,;\,\frac{1\,140}{3}\right) = (5\,;\,380)$.

\emph{Équation.} On cherche $a$ et $b$ tels que
\[
  \begin{cases} 2a+b = 325\\ 5a+b = 380 \end{cases}
\]
En soustrayant : $3a = 55$, donc $a = \frac{55}{3} \approx 18{,}33$ et
$b = 325-\frac{110}{3} = \frac{865}{3} \approx 288{,}33$. La droite a pour
équation
\[
  y = \frac{55}{3}\,x+\frac{865}{3} .
\]

\emph{Vérification.} Pour $x = 3{,}5$ :
$\frac{55\times3{,}5+865}{3} = \frac{1\,057{,}5}{3} = 352{,}5 = \moy{y}$.
\checkmark

\emph{Estimation.} L'année $2030$ porte le rang $12$ :
$y = \frac{55\times12+865}{3} = \frac{1\,525}{3} \approx 508$. On peut
estimer à environ $508$ le nombre d'inscrits en $2030$.
\end{exemple}

\begin{avertissement}
Une estimation faite à partir d'une droite d'ajustement suppose que la
tendance observée se poursuit. Rien ne le garantit — une épidémie, une
réforme, une sécheresse changent la donne. On écrit donc « on peut
\emph{estimer} », jamais « il y aura ».
\end{avertissement}

\begin{entrainement}[nuages et ajustements]
\exo[\niveau{1}] Calculer le point moyen de la série
$(1\,;\,4)$, $(2\,;\,7)$, $(3\,;\,9)$, $(4\,;\,12)$.

\exo[\niveau{2}] Pour cette même série, déterminer $G_{1}$ et $G_{2}$, puis
l'équation de $\droiteMayer$.

\exo[\niveau{2}] Vérifier que le point moyen appartient à cette droite.

\exo[\niveau{3}] Une série de six couples a pour points moyens partiels
$G_{1}(2\,;\,66)$ et $G_{2}(5\,;\,75)$. Déterminer l'équation de la droite
de Mayer, puis estimer $y$ pour $x = 11$.

\exo[\niveau{3}] Pour la droite $y = 1{,}0625\,x+2{,}0375$, estimer la
valeur de $x$ correspondant à $y = 10{,}5$.
\end{entrainement}

\begin{corrige}[nuages et ajustements]
\textbf{1.} $\moy{x} = 2{,}5$ et $\moy{y} = 8$ : $\pmoyen(2{,}5\,;\,8)$.
\textbf{2.} $G_{1}(1{,}5\,;\,5{,}5)$ et $G_{2}(3{,}5\,;\,10{,}5)$ ; la
pente vaut $\frac{5}{2} = 2{,}5$ et $b = 5{,}5-3{,}75 = 1{,}75$ :
$y = 2{,}5x+1{,}75$.
\textbf{3.} Pour $x = 2{,}5$ : $6{,}25+1{,}75 = 8 = \moy{y}$. \checkmark
\textbf{4.} $3a = 9$ donc $a = 3$ et $b = 60$ : $y = 3x+60$ ; pour
$x = 11$, $y = 93$.
\textbf{5.} $1{,}0625x = 8{,}4625$, donc $x \approx 7{,}96$, soit environ
$8$.
\end{corrige}

% ===========================================================================
\section{Taux d'évolution}

\begin{definition}[taux d'évolution]
Lorsqu'une grandeur passe de la valeur $V_{0}$ à la valeur $V_{1}$, son
\textbf{taux d'évolution} est
\[
  t = \frac{V_{1}-V_{0}}{V_{0}},
\]
que l'on exprime en pourcentage. Le \textbf{coefficient multiplicateur}
associé vaut $1+t$.
\end{definition}

\begin{propriete}[évolutions successives]
Deux évolutions successives de taux $t_{1}$ puis $t_{2}$ équivalent à une
évolution unique de coefficient multiplicateur
$\left(1+t_{1}\right)\left(1+t_{2}\right)$. Les taux, eux, ne s'additionnent
pas.
\end{propriete}

\begin{erreurclassique}
Une hausse de $20\,\%$ suivie d'une baisse de $20\,\%$ ne ramène pas au
point de départ : le coefficient global vaut
$1{,}2\times0{,}8 = 0{,}96$, soit une baisse de $4\,\%$. C'est l'exemple à
retenir, et à ressortir chaque fois que l'on est tenté d'additionner deux
pourcentages.
\end{erreurclassique}

\begin{entrainement}[pourcentages]
\exo[\niveau{1}] Une production passe de $80$ à $92$ tonnes : calculer le
taux d'évolution.

\exo[\niveau{2}] Un prix augmente de $15\,\%$ puis de $10\,\%$ : quel est le
taux global ?

\exo[\niveau{2}] Un effectif diminue de $8\,\%$ puis augmente de $8\,\%$ :
quel est le taux global ?

\exo[\niveau{3}] Un capital augmente de $5\,\%$ par an. Au bout de combien
d'années aura-t-il doublé ? On donne $\ln 2 \approx 0{,}693$ et
$\ln 1{,}05 \approx 0{,}0488$.
\end{entrainement}

\begin{corrige}[pourcentages]
\textbf{1.} $\frac{92-80}{80} = 0{,}15$, soit $+15\,\%$.
\textbf{2.} $1{,}15\times1{,}10 = 1{,}265$, soit $+26{,}5\,\%$.
\textbf{3.} $0{,}92\times1{,}08 = 0{,}9936$, soit $-0{,}64\,\%$.
\textbf{4.} On cherche $n$ tel que $1{,}05^{n} = 2$, soit
$n = \frac{\ln2}{\ln1{,}05} \approx 14{,}2$ : au bout de quinze ans.
\end{corrige}

\begin{qcm}
\exo Le point moyen d'une série de six couples se calcule avec :
\textbf{a)} les trois premiers couples \quad \textbf{b)} tous les couples
\quad \textbf{c)} les valeurs extrêmes \quad \textbf{d)} les $x_{i}$
seulement \reponseqcm{b}

\exo La droite de Mayer passe nécessairement par :
\textbf{a)} l'origine \quad \textbf{b)} le premier point du nuage \quad
\textbf{c)} le point moyen $\pmoyen$ \quad \textbf{d)} tous les points du
nuage \reponseqcm{c}

\exo $G_{1}(2\,;\,34)$ et $G_{2}(5\,;\,44)$ : le coefficient directeur de
$\droiteMayer$ vaut :
\textbf{a)} $10$ \quad \textbf{b)} $\frac{10}{3}$ \quad \textbf{c)} $3$
\quad \textbf{d)} $\frac{3}{10}$ \reponseqcm{b}

\exo Une hausse de $10\,\%$ correspond au coefficient multiplicateur :
\textbf{a)} $0{,}10$ \quad \textbf{b)} $1{,}10$ \quad \textbf{c)} $10$
\quad \textbf{d)} $0{,}90$ \reponseqcm{b}

\exo Estimer une valeur au-delà des données observées s'appelle :
\textbf{a)} une interpolation \quad \textbf{b)} une extrapolation \quad
\textbf{c)} une moyenne \quad \textbf{d)} un ajustement \reponseqcm{b}
\end{qcm}

% ===========================================================================
\begin{annales}
La statistique double occupe l'un des deux exercices de cinq points à peu
près une session sur deux. Les énoncés ci-dessous sont repris tels quels.

\exoann{Baccalauréat, Madagascar, série A, session 2026 — exercice 2}
Le tableau ci-dessous donne l'évolution de l'effectif des élèves inscrits au
baccalauréat durant six années consécutives.

\begin{center}\footnotesize
\renewcommand{\arraystretch}{1.25}
\begin{tabular}{|l|c|c|c|c|c|c|}
\hline
Année & $2019$ & $2020$ & $2021$ & $2022$ & $2023$ & $2024$\\
\hline
Rang de l'année $x_{i}$ & $1$ & $2$ & $3$ & $4$ & $5$ & $6$\\
\hline
Nombre d'élèves $y_{i}$ & $340$ & $310$ & $325$ & $380$ & $365$ & $395$\\
\hline
\end{tabular}
\end{center}

\begin{enumerate}[label=\textbf{\arabic*)}, leftmargin=*, itemsep=0.15em]
  \item Construire le nuage de points associé à cette série statistique dans
        un repère orthogonal. Échelle : sur l'axe des abscisses, $1$ cm pour
        $1$ unité ; sur l'axe des ordonnées, placer $300$ à l'origine puis
        $1$ cm pour $10$ élèves.
  \item Déterminer les coordonnées du point moyen $G$.
  \item On partage cette série en deux sous-groupes d'effectifs égaux, le
        premier composé des trois premières années et le second des trois
        dernières.
        \begin{enumerate}[label=\textbf{\alph*)}, leftmargin=1.4em]
          \item Calculer les coordonnées respectives des points moyens
                $G_{1}$ et $G_{2}$ de ces deux sous-groupes.
          \item Établir l'équation de la droite d'ajustement
                $\left(G_{1}G_{2}\right)$ par la méthode de Mayer.
        \end{enumerate}
  \item En utilisant cette droite, estimer l'effectif des élèves inscrits au
        baccalauréat en $2030$.
\end{enumerate}

\exoann{Baccalauréat, Madagascar, série A, session 2017 — exercice 2}
Le tableau suivant représente l'évolution des bénéfices par mois d'un
marchand ($y_{i}$, en millions d'ariary), où $y_{6}$ est un entier naturel.

\begin{center}\footnotesize
\renewcommand{\arraystretch}{1.25}
\begin{tabular}{|l|c|c|c|c|c|c|}
\hline
Mois & Janvier & Février & Mars & Avril & Mai & Juin\\
\hline
Rang du mois $x_{i}$ & $1$ & $2$ & $3$ & $4$ & $5$ & $6$\\
\hline
Bénéfice $y_{i}$ & $32$ & $34$ & $36$ & $41$ & $42$ & $y_{6}$\\
\hline
\end{tabular}
\end{center}

\begin{enumerate}[label=\textbf{\arabic*)}, leftmargin=*, itemsep=0.15em]
  \item Calculer la valeur de $y_{6}$ si la moyenne de la série $(y_{i})$
        est $\moy{y} = 39$.
  \item Dans tout ce qui suit, on prendra $y_{6} = 49$.
        \begin{enumerate}[label=\textbf{\alph*)}, leftmargin=1.4em]
          \item Représenter le nuage de points $M_{i}(x_{i}\,;\,y_{i})$ dans
                un repère orthogonal. Sur l'axe des abscisses, prendre
                $1$ cm comme unité ; sur l'axe des ordonnées, placer $30$ à
                l'origine et prendre $1$ cm pour représenter $2$ millions
                d'ariary.
          \item Déterminer les coordonnées du point moyen $G$.
          \item Écrire l'équation cartésienne de la droite d'ajustement
                linéaire $\left(G_{1}G_{2}\right)$ par la méthode de Mayer.
        \end{enumerate}
  \item À l'aide de cette droite, estimer le bénéfice du marchand au mois de
        septembre.
\end{enumerate}

\exoann{Baccalauréat, Madagascar, série A, session 2014 — exercice 2}
L'évolution de la production rizicole $y_{i}$, en tonnes, d'un Petit
Périmètre Irrigué durant six années est représentée par le tableau suivant.

\begin{center}\footnotesize
\renewcommand{\arraystretch}{1.25}
\begin{tabular}{|l|c|c|c|c|c|c|}
\hline
Année & $2008$ & $2009$ & $2010$ & $2011$ & $2012$ & $2013$\\
\hline
Rang $x_{i}$ & $1$ & $2$ & $3$ & $4$ & $5$ & $6$\\
\hline
Production $y_{i}$ & $65$ & $70$ & $63$ & $68$ & $76$ & $81$\\
\hline
\end{tabular}
\end{center}

\begin{enumerate}[label=\textbf{\arabic*)}, leftmargin=*, itemsep=0.15em]
  \item Construire le nuage de points associé. Sur l'axe des abscisses,
        $1$ cm pour une unité ; sur l'axe des ordonnées, $60$ à l'origine et
        $1$ cm pour $5$ tonnes.
  \item Calculer les coordonnées du point moyen $G$.
  \item En utilisant la méthode de Mayer, déterminer une équation de la
        droite d'ajustement.
  \item Tracer cette droite.
  \item Déterminer graphiquement la production du périmètre en $2018$, puis
        vérifier le résultat par le calcul.
\end{enumerate}

\exoann{Baccalauréat, Madagascar, série A, session 2010 — exercice 1}
Les notes obtenues par dix candidats aux épreuves de natation et de danse
d'un concours sont indiquées dans le tableau suivant.

\begin{center}\footnotesize
\renewcommand{\arraystretch}{1.25}
\begin{tabular}{|l|c|c|c|c|c|c|c|c|c|c|}
\hline
Natation $(X)$ & $3$ & $5$ & $6$ & $6$ & $9$ & $9$ & $12$ & $12$ & $14$ &
$14$\\
\hline
Danse $(Y)$ & $5$ & $8$ & $8$ & $10$ & $10$ & $13$ & $13$ & $16$ & $16$ &
$17$\\
\hline
\end{tabular}
\end{center}

\begin{enumerate}[label=\textbf{\arabic*)}, leftmargin=*, itemsep=0.15em]
  \item Représenter le nuage de points associé à cette série statistique.
  \item On partage le nuage en deux parties d'effectifs égaux.
        \begin{enumerate}[label=\textbf{\alph*)}, leftmargin=1.4em]
          \item Calculer les coordonnées de $G_{1}$ et $G_{2}$, points
                moyens respectifs des nuages partiels ainsi obtenus.
          \item Placer $G_{1}$ et $G_{2}$, puis tracer la droite
                $\left(G_{1}G_{2}\right)$.
          \item Déterminer une équation de la droite
                $\left(G_{1}G_{2}\right)$.
        \end{enumerate}
  \item Estimer, à l'aide de cette droite, la note de natation qu'aura un
        candidat qui avait $10{,}5$ en danse.
\end{enumerate}

\exoann{Baccalauréat, Madagascar, série A, session 2002 — exercice 2}
Le tableau suivant indique l'évolution de l'effectif d'un collège au cours
des huit dernières années, $x_{i}$ désignant le rang de l'année et $y_{i}$
l'effectif correspondant : pour les rangs $1$ à $8$ (années $1994$ à
$2001$), les effectifs sont $370$, $360$, $380$, $410$, $420$, $440$, $450$
et $470$.

\begin{enumerate}[label=\textbf{\arabic*)}, leftmargin=*, itemsep=0.15em]
  \item Représenter le nuage de points dans un repère orthogonal. Sur l'axe
        des abscisses, $1$ cm pour unité ; sur l'axe des ordonnées, placer
        $350$ à l'origine puis $1$ cm pour $10$ élèves.
  \item Calculer les coordonnées du point moyen $G$.
  \item On note $(S_{1})$ la série allant de $1994$ à $1997$ et $(S_{2})$
        celle allant de $1998$ à $2001$.
        \begin{enumerate}[label=\textbf{\alph*)}, leftmargin=1.4em]
          \item Déterminer les coordonnées des points moyens $G_{1}$ et
                $G_{2}$.
          \item Déterminer une équation cartésienne de la droite $(D)$
                passant par $G_{1}$ et $G_{2}$.
          \item Construire $(D)$. Que représente-t-elle ?
          \item En déduire une estimation de l'effectif du collège en
                $2003$.
        \end{enumerate}
\end{enumerate}

\exoann{Baccalauréat, Madagascar, série A, session 2000 — exercice 2}
Le tableau suivant montre le chiffre d'affaires, en millions de francs
malagasy, d'une entreprise au cours de six années : pour les rangs $1$ à $6$
(années $1994$ à $1999$), les chiffres d'affaires sont $120$, $132$, $147$,
$164$, $181$ et $201$.

\begin{enumerate}[label=\textbf{\arabic*)}, leftmargin=*, itemsep=0.15em]
  \item Calculer la moyenne de la série $\left(y_{i}\right)$.
  \item Représenter le nuage de points $M_{i}\left(x_{i}\,;\,y_{i}\right)$.
  \item Soit $G_{1}$ le point moyen du sous-nuage des trois premières
        années et $G_{2}$ celui des trois dernières.
        \begin{enumerate}[label=\textbf{\alph*)}, leftmargin=1.4em]
          \item Déterminer les coordonnées de $G_{1}$ et $G_{2}$.
          \item Tracer la droite $\left(G_{1}G_{2}\right)$ ; que
                représente-t-elle ?
          \item Donner l'équation de cette droite.
          \item En déduire une prévision du chiffre d'affaires en $2002$.
        \end{enumerate}
\end{enumerate}
\end{annales}

\begin{corrige}[annales]
\textbf{1. (2026)} Voir l'exemple du cours : $G(3{,}5\,;\,352{,}5)$,
$G_{1}(2\,;\,325)$, $G_{2}(5\,;\,380)$,
$y = \frac{55}{3}x+\frac{865}{3} \approx 18{,}33x+288{,}33$, et
l'estimation pour $2030$ (rang $12$) vaut environ $508$ élèves.

\textbf{2. (2017)} La somme des six bénéfices vaut $6\times39 = 234$, et
$32+34+36+41+42 = 185$ : donc $y_{6} = 49$. Ensuite
$G(3{,}5\,;\,39)$, $G_{1}\left(2\,;\,34\right)$ et
$G_{2}\left(5\,;\,44\right)$. La pente vaut $\frac{44-34}{3} = \frac{10}{3}$
et $b = 34-\frac{20}{3} = \frac{82}{3}$, d'où
$y = \frac{10}{3}x+\frac{82}{3}$. Septembre porte le rang $9$ :
$y = \frac{90+82}{3} = \frac{172}{3} \approx 57{,}3$ millions d'ariary.

\textbf{3. (2014)} $\moy{y} = \frac{423}{6} = 70{,}5$, donc
$G(3{,}5\,;\,70{,}5)$. $G_{1}(2\,;\,66)$ et $G_{2}(5\,;\,75)$ : la pente
vaut $3$ et $b = 60$, d'où $y = 3x+60$. En $2018$, de rang $11$ :
$y = 93$ tonnes.

\textbf{4. (2010)} $G_{1}$ résume les cinq premiers candidats :
$\left(\frac{29}{5}\,;\,\frac{41}{5}\right) = (5{,}8\,;\,8{,}2)$ ; $G_{2}$
les cinq derniers : $\left(\frac{61}{5}\,;\,15\right) = (12{,}2\,;\,15)$.
La pente vaut $\frac{15-8{,}2}{12{,}2-5{,}8} = \frac{6{,}8}{6{,}4}
= 1{,}0625$ et $b = 8{,}2-1{,}0625\times5{,}8 = 2{,}0375$, d'où
$y = 1{,}0625x+2{,}0375$. Pour $y = 10{,}5$ :
$x = \frac{10{,}5-2{,}0375}{1{,}0625} \approx 7{,}96$, soit environ $8$ en
natation.

\textbf{5. (2002)} $\moy{x} = 4{,}5$ et
$\moy{y} = \frac{3\,300}{8} = 412{,}5$. $G_{1}(2{,}5\,;\,380)$ et
$G_{2}(6{,}5\,;\,445)$ : la pente vaut $\frac{65}{4} = 16{,}25$ et
$b = 380-40{,}625 = 339{,}375$, d'où $y = 16{,}25x+339{,}375$. Cette droite
est la droite d'ajustement de la série. En $2003$, de rang $10$ :
$y \approx 502$ élèves.

\textbf{6. (2000)} $\moy{y} = \frac{945}{6} = 157{,}5$.
$G_{1}(2\,;\,133)$ et $G_{2}(5\,;\,182)$ : la pente vaut
$\frac{49}{3} \approx 16{,}33$ et
$b = 133-\frac{98}{3} = \frac{301}{3} \approx 100{,}33$, d'où
$y = \frac{49}{3}x+\frac{301}{3}$. En $2002$, de rang $9$ :
$y = \frac{441+301}{3} \approx 247$ millions de francs malagasy.
\end{corrige}

% ===========================================================================
\begin{vraifaux}
\exo Le point moyen appartient toujours au nuage de points.

\exo La droite de Mayer passe par $G_{1}$ et $G_{2}$.

\exo Deux hausses successives de $50\,\%$ équivalent à une hausse de
$100\,\%$.

\exo Une estimation par extrapolation est une certitude.

\exo Si l'on échange les rôles de $x$ et de $y$, on obtient la même droite.
\end{vraifaux}

\begin{corrige}[vrai ou faux]
\textbf{1.} Faux — il est au centre du nuage mais n'en fait pas
nécessairement partie ; dans l'exemple du cours, $G(3{,}5\,;\,352{,}5)$
n'est aucun des six points.
\textbf{2.} Vrai — c'est sa définition même.
\textbf{3.} Faux — le coefficient global vaut $1{,}5\times1{,}5 = 2{,}25$,
soit une hausse de $125\,\%$.
\textbf{4.} Faux — elle suppose que la tendance se poursuit.
\textbf{5.} Faux — on obtient une autre droite, qui sert à estimer $x$
connaissant $y$ ; c'est d'ailleurs ce que fait la dernière question du
sujet 2010, en résolvant l'équation plutôt qu'en changeant de droite.
\end{corrige}

% ===========================================================================
\begin{situation}[le prix du riz au marché]
Un observateur relève, le premier lundi de chaque mois, le prix du kilo de
riz sur un marché d'Antananarivo, en ariary.

\begin{center}\footnotesize
\renewcommand{\arraystretch}{1.25}
\begin{tabular}{|l|c|c|c|c|c|c|}
\hline
Rang du mois $x_{i}$ & $1$ & $2$ & $3$ & $4$ & $5$ & $6$\\
\hline
Prix $y_{i}$ & $2\,800$ & $2\,900$ & $3\,050$ & $3\,100$ & $3\,250$ &
$3\,300$\\
\hline
\end{tabular}
\end{center}

\begin{enumerate}[label=\textbf{\arabic*.}, leftmargin=*, itemsep=0.3em]
  \item Calculer le point moyen $\pmoyen$ de la série.
  \item Déterminer $G_{1}$ et $G_{2}$, puis l'équation de la droite
        d'ajustement par la méthode de Mayer.
  \item Vérifier que $\pmoyen$ appartient bien à cette droite.
  \item Estimer le prix du kilo de riz au dixième mois.
  \item Calculer le taux d'évolution du prix entre le premier et le sixième
        mois, puis dire si la droite trouvée en rend compte fidèlement.
\end{enumerate}
\end{situation}

\begin{corrige}[le prix du riz au marché]
\textbf{1.} $\moy{x} = 3{,}5$ et
$\moy{y} = \frac{18\,400}{6} \approx 3\,066{,}67$.

\textbf{2.} $G_{1}\left(2\,;\,\frac{8\,750}{3}\right)
\approx (2\,;\,2\,916{,}67)$ et
$G_{2}\left(5\,;\,\frac{9\,650}{3}\right) \approx (5\,;\,3\,216{,}67)$. La
pente vaut $\frac{900}{9} = 100$ et
$b = 2\,916{,}67-200 = 2\,716{,}67$, d'où
$y = 100x+2\,716{,}67$.

\textbf{3.} Pour $x = 3{,}5$ : $350+2\,716{,}67 = 3\,066{,}67 = \moy{y}$.
\checkmark

\textbf{4.} $y = 1\,000+2\,716{,}67 \approx 3\,717$ ariary.

\textbf{5.} $\frac{3\,300-2\,800}{2\,800} \approx 0{,}179$, soit environ
$+17{,}9\,\%$ en cinq mois. La droite, qui prévoit une hausse de $100$
ariary par mois, en rend bien compte : $500$ ariary sur cinq mois, ce qui
correspond exactement à la hausse observée.
\end{corrige}

% ===========================================================================
\begin{jemeteste}
Vingt-cinq minutes.

\begin{enumerate}[label=\textbf{\arabic*.}, leftmargin=*, itemsep=0.3em]
  \item Une série a pour valeurs $x$ : $1$, $2$, $3$, $4$, $5$, $6$ et pour
        valeurs $y$ : $12$, $15$, $14$, $19$, $21$, $23$. Calculer
        $\pmoyen$.
  \item Déterminer $G_{1}$ et $G_{2}$.
  \item Établir l'équation de la droite de Mayer.
  \item Vérifier que $\pmoyen$ y appartient.
  \item Estimer $y$ pour $x = 10$.
\end{enumerate}
\end{jemeteste}

\begin{corrige}[je me teste]
\textbf{1.} $\moy{x} = 3{,}5$ et $\moy{y} = \frac{104}{6}
\approx 17{,}33$.
\textbf{2.} $G_{1}\left(2\,;\,\frac{41}{3}\right) \approx (2\,;\,13{,}67)$
et $G_{2}\left(5\,;\,21\right)$.
\textbf{3.} La pente vaut $\frac{21-13{,}67}{3} \approx 2{,}44$ et
$b \approx 13{,}67-4{,}89 = 8{,}78$ : $y \approx 2{,}44x+8{,}78$.
\textbf{4.} Pour $x = 3{,}5$ : $8{,}55+8{,}78 = 17{,}33$. \checkmark
\textbf{5.} $y \approx 24{,}4+8{,}78 \approx 33{,}2$.
\end{corrige}

% ===========================================================================
\begin{commeaubac}
\bmtsource{Baccalauréat, Madagascar, série A, session 2026 — exercice 2
\textbullet\ 5 points}
Le tableau ci-dessous consigne l'évolution de l'effectif des élèves inscrits
à l'examen du baccalauréat durant six années consécutives.

\begin{center}\footnotesize
\renewcommand{\arraystretch}{1.25}
\begin{tabular}{|l|c|c|c|c|c|c|}
\hline
Année & $2019$ & $2020$ & $2021$ & $2022$ & $2023$ & $2024$\\
\hline
Rang de l'année $x_{i}$ & $1$ & $2$ & $3$ & $4$ & $5$ & $6$\\
\hline
Nombre d'élèves $y_{i}$ & $340$ & $310$ & $325$ & $380$ & $365$ & $395$\\
\hline
\end{tabular}
\end{center}

\begin{enumerate}[label=\textbf{\arabic*.}, leftmargin=*, itemsep=0.25em]
  \item Construire le nuage des points associé à cette série statistique
        dans un repère orthogonal d'échelle : sur l'axe des abscisses,
        $1$ cm pour $1$ unité ; sur l'axe des ordonnées, placer $300$ à
        l'origine puis $1$ cm pour $10$ élèves. \bareme{1 ; 1}
  \item Déterminer les coordonnées du point moyen $G$. \bareme{1 ; 1}
  \item On partage cette série en deux sous-groupes d'effectifs égaux. La
        première partie est composée des trois premières années et la
        seconde des trois dernières.
    \begin{enumerate}[label=\textbf{\alph*)}, leftmargin=1.4em, itemsep=0.15em]
      \item Calculer les coordonnées respectives des points moyens $G_{1}$
            et $G_{2}$ de ces deux sous-groupes. \bareme{1 ; 1}
      \item Établir l'équation de la droite d'ajustement
            $\left(G_{1}G_{2}\right)$ par la méthode de Mayer.
            \bareme{1 ; 1}
    \end{enumerate}
  \item En utilisant cette droite, estimer l'effectif des élèves inscrits au
        baccalauréat en $2030$. \bareme{1 ; 1}
\end{enumerate}
\end{commeaubac}

\begin{corrige}[comme au baccalauréat — Madagascar A 2026]
\textbf{1.} On place les six points $(1\,;\,340)$, $(2\,;\,310)$,
$(3\,;\,325)$, $(4\,;\,380)$, $(5\,;\,365)$ et $(6\,;\,395)$ dans le repère
demandé, en marquant la rupture d'échelle sur l'axe des ordonnées juste
au-dessus de $300$.

\textbf{2.} $\moy{x} = \frac{1+2+3+4+5+6}{6} = 3{,}5$ et
$\moy{y} = \frac{340+310+325+380+365+395}{6} = \frac{2\,115}{6} = 352{,}5$.
Donc $G(3{,}5\,;\,352{,}5)$.

\textbf{3. a)} $G_{1}\left(\frac{1+2+3}{3}\,;\,\frac{340+310+325}{3}\right)
= (2\,;\,325)$ et
$G_{2}\left(\frac{4+5+6}{3}\,;\,\frac{380+365+395}{3}\right)
= (5\,;\,380)$.
\textbf{b)} La droite $\left(G_{1}G_{2}\right)$ a pour équation
$y = ax+b$ avec
\[
  \begin{cases} 2a+b = 325\\ 5a+b = 380 \end{cases}
  \quad\Longrightarrow\quad
  3a = 55, \quad a = \frac{55}{3}, \quad b = \frac{865}{3}.
\]
Soit $y = \frac{55}{3}x+\frac{865}{3}$, c'est-à-dire environ
$y = 18{,}33\,x+288{,}33$.
On vérifie que $G$ y appartient : pour $x = 3{,}5$, on retrouve $352{,}5$.

\textbf{4.} L'année $2030$ porte le rang $x = 12$ :
\[
  y = \frac{55\times12+865}{3} = \frac{1\,525}{3} \approx 508{,}33 .
\]
On peut estimer à environ $508$ le nombre d'élèves inscrits en $2030$.
\end{corrige}

\begin{notepourlaclasse}
Huit heures, et c'est le chapitre le plus rentable de la partie : les
techniques sont peu nombreuses et la question tombe presque à chaque
session.

Le tracé compte pour un point entier. Exigez le papier millimétré, la
rupture d'échelle correctement marquée sur l'axe des ordonnées, et les
points bien visibles. Un nuage tracé au jugé, sans respecter l'échelle
imposée, ne rapporte rien.

Le calcul de $G_{1}$ et $G_{2}$ se fait presque toujours sur des rangs
consécutifs : $G_{1}$ a pour abscisse $2$ et $G_{2}$ pour abscisse $5$
lorsque la série compte six valeurs. Faites-le remarquer, mais faites
recalculer : deux sessions sur les six du bloc d'annales utilisent d'autres
effectifs.

La vérification « $G$ appartient à la droite » n'est jamais demandée par les
sujets, mais elle détecte la quasi-totalité des erreurs de calcul. Imposez-la
comme un réflexe.
\end{notepourlaclasse}
